% Separate notations: dashed enclosures = e-classes; solid small boxes = e-nodes.
% Diagram arrows are explicitly labelled when they denote execution or catalog
% references rather than child slots. No colour memorisation is required.
\tikzset{en/.style={draw,rounded corners=1pt,fill=white,inner sep=3pt,font=\small},
 ec/.style={draw,dashed,rounded corners,inner sep=5pt},
 note/.style={align=center,font=\small},
 rec/.style={draw,align=center,text width=43mm,inner sep=5pt,font=\small}}
\newcommand{\TeachingReplay}{%
\begin{figure}[!htbp]\centering
\begin{tikzpicture}
\node[note] at (0,0.5) {\textbf{(a) Entry $T_R$}};
\node[note] at (8,0.5) {\textbf{(b) Entry $T_L$}};
\node[en] (r) at (0,0) {$a+(b+c)$};
\node[en] (l) at (8,0) {$(a+b)+c$};
\node[en] (r1) at (0,-1.5) {$a+(b+c)$};
\node[en] (r2) at (0,-2.15) {$(a+b)+c$};
\node[en] (l1) at (8,-1.5) {$(a+b)+c$};
\node[en] (l2) at (8,-2.15) {$c+(a+b)$};
\node[ec,fit=(r1)(r2)] (rc) {};
\node[ec,fit=(l1)(l2)] (lc) {};
\draw[arr] (r)--node[right,font=\scriptsize]{assoc + union}(rc);
\draw[arr] (l)--node[right,font=\scriptsize]{comm + union}(lc);
\node[note,text width=125mm] (end) at (4,-3.7)
 {\textbf{Continue both rules to a fixed point}\\
 Each root class has six alternatives; each whole Math graph has\\
 12 Add nodes + 3 Var nodes in 7 e-classes (Figure~\ref{fig:body}).};
\draw[arr] (rc.south)--(end.north west);
\draw[arr] (lc.south)--(end.north east);
\end{tikzpicture}
\caption{Two familiar expressions, two local executions. Dashed enclosures show
one root e-class with retained alternatives. Expressions abbreviate e-node graphs;
child classes are expanded in Figure~\ref{fig:body}. Labelled arrows denote
selected rewrite steps, not deletion or replacement of the original expression.
The intermediate panels explain individual applications; the automated ripen test
checks the final fixed points rather than prescribing these first applications.}
\label{fig:replay}
\end{figure}}

\newcommand{\TeachingCatalog}{%
\begin{figure}[!htbp]\centering
\begin{tikzpicture}
\node[rec] (r) at (0,0) {$T_R$: $a+(b+c)$\\own entry and evidence};
\node[rec] (l) at (8,0) {$T_L$: $(a+b)+c$\\own entry and evidence};
\node[rec] (kr) at (0,-1.6) {$K_R$: 12 Add + 3 Var\\7 Math e-classes};
\node[rec] (kl) at (8,-1.6) {$K_L$: 12 Add + 3 Var\\7 Math e-classes};
\draw[arr] (r)--node[right,font=\scriptsize]{native ripen}(kr);
\draw[arr] (l)--node[right,font=\scriptsize]{native ripen}(kl);
\node[note] at (4,-2.65) {\textbf{Exact graph comparison succeeds}};
\node[rec] (tr) at (0,-4) {Catalog trigger $T_R$\\entry + value map};
\node[rec] (tl) at (8,-4) {Catalog trigger $T_L$\\entry + value map};
\node[rec,text width=35mm] (k) at (4,-5.5) {One representative $K$\\same rows and ports};
\draw[arr] (tr)--node[left,font=\scriptsize]{references}(k);
\draw[arr] (tl)--node[right,font=\scriptsize]{references}(k);
\draw[dotted] (-2.5,-3.1)--(10.5,-3.1);
\end{tikzpicture}
\caption{Before comparison: two completed replay exports. After comparison:
two catalog trigger records reference one representative body. Solid rectangles
are records, not e-classes. Equal counts alone do not justify either reference:
the checker requires a complete typed bijection preserving roots, rows, and
visibility. Per-run exports and the live tier0 graph are not deleted.}
\label{fig:catalog}
\end{figure}}

\newcommand{\TeachingBody}{%
\begin{figure}[!htbp]\centering
\begin{tikzpicture}
\node[en] (r0) at (0,0) {$\mathrm{Add}(a,q_{bc})$};
\node[en] (r1) at (4,0) {$\mathrm{Add}(q_{bc},a)$};
\node[en] (r2) at (8,0) {$\mathrm{Add}(b,q_{ac})$};
\node[en] (r3) at (0,-0.75) {$\mathrm{Add}(q_{ac},b)$};
\node[en] (r4) at (4,-0.75) {$\mathrm{Add}(c,q_{ab})$};
\node[en] (r5) at (8,-0.75) {$\mathrm{Add}(q_{ab},c)$};
\node[ec,fit=(r0)(r1)(r2)(r3)(r4)(r5),label=above:{root e-class $r$}] {};
\node[en] (ab0) at (0,-2.3) {$\mathrm{Add}(a,b)$};
\node[en] (ab1) at (0,-2.95) {$\mathrm{Add}(b,a)$};
\node[ec,fit=(ab0)(ab1),label=above:{$q_{ab}$}] (ab) {};
\node[en] (ac0) at (4,-2.3) {$\mathrm{Add}(a,c)$};
\node[en] (ac1) at (4,-2.95) {$\mathrm{Add}(c,a)$};
\node[ec,fit=(ac0)(ac1),label=above:{$q_{ac}$}] {};
\node[en] (bc0) at (8,-2.3) {$\mathrm{Add}(b,c)$};
\node[en] (bc1) at (8,-2.95) {$\mathrm{Add}(c,b)$};
\node[ec,fit=(bc0)(bc1),label=above:{$q_{bc}$}] {};
\node[en] (a0) at (0,-4.4) {Var("a")};
\node[en] (b0) at (4,-4.4) {Var("b")};
\node[en] (c0) at (8,-4.4) {Var("c")};
\node[ec,fit=(a0),label=below:{$a$}] (a) {};
\node[ec,fit=(b0),label=below:{$b$}] (b) {};
\node[ec,fit=(c0),label=below:{$c$}] {};
\draw[arr] (ab1.south)--node[left,font=\scriptsize]{slot 1}(a.north);
\draw[arr] (ab1.south east)--node[above,font=\scriptsize]{slot 0}(b.north west);
\end{tikzpicture}
\caption{Count and read the entire saturated Math body. Each solid box is one
e-node; each dashed enclosure is one e-class. The ordered labels specify all
child-class references. Two arrows expand the slots of Add$(b,a)$ as a reading
example; the other references remain explicit in the labels. There are six root
Add nodes, six pair Add nodes, and three Var nodes. Primitive String values are
omitted. Only association and commutation are enabled.}
\label{fig:body}
\end{figure}}

\newcommand{\TeachingBoundary}{%
\begin{figure}[!htbp]\centering
\begin{tikzpicture}
\node[rec] (i) at (0,0) {External input\\$a+(b+c)$};
\node[rec] (a) at (6,0) {Application 1: assoc\\produces $(a+b)+c$};
\node[rec] (b) at (6,-1.7) {Application 2: comm\\produces $c+(a+b)$};
\draw[arr] (i)--node[above,font=\scriptsize]{supplies premise}(a);
\draw[arr] (a)--node[right,font=\scriptsize]{produced row is read}(b);
\node[draw,dashed,fit=(a)(b),inner sep=7pt,label=below:{selected two-application boundary $H$}] {};
\node[note,text width=42mm] at (0,-1.7) {Inside $H$: 1 is coarse,\\2 is smooth.\\Alone: 2 needs an external row.};
\end{tikzpicture}
\caption{The boundary changes the role, not the rule. Rectangles represent
application or input records; the dashed enclosure is the selected history
fragment, not an e-class. The dependency arrow identifies a producer-to-consumer
row witness. Moving application 1 outside the boundary makes that dependency
external for application 2.}
\label{fig:boundary}
\end{figure}}
